
\documentclass{article}
\usepackage{colortbl}
\usepackage{makecell}
\usepackage{multirow}
\usepackage{supertabular}

\begin{document}

\newcounter{utterance}

\centering \large Interaction Transcript for game `wordle', experiment `medium\_frequency\_words\_no\_clue\_no\_critic', episode 3 with gemma{-}4{-}26b{-}a4b{-}it.
\vspace{24pt}

{ \footnotesize  \setcounter{utterance}{1}
\setlength{\tabcolsep}{0pt}
\begin{supertabular}{c@{$\;$}|p{.15\linewidth}@{}p{.15\linewidth}p{.15\linewidth}p{.15\linewidth}p{.15\linewidth}p{.15\linewidth}}
    \# & \multicolumn{2}{c}{Player} && \multicolumn{2}{c}{Game Master} \\
    \hline

    \theutterance \stepcounter{utterance}  
    & & & \multicolumn{4}{p{0.6\linewidth}}{
        \cellcolor[rgb]{0.9,0.9,0.9}{
            \makecell[{{p{\linewidth}}}]{
                \texttt{\tiny{[P1$\langle$GM]}}
                \texttt{Jste jazykový kouzelník, který rád hádá slova podle daných pravidel.} \\
\\ 
\texttt{Vítejte ve hře Wordle! Máte šest pokusů na uhodnutí cílového slova, platného českého slova o pěti malých písmenech. Použijte značky "vysvětlení:" a "hádání:" k poskytnutí stručného vysvětlení pro každý odhad.} \\
\\ 
\texttt{Například, pokud je váš odhad "jablko", vaše odpověď by měla být} \\
\texttt{vysvětlení: toto je běžné české slovo o pěti písmenech a začínám svůj odhad tímto slovem.} \\
\texttt{hádání: jablko} \\
\\ 
\texttt{Po každém odhadu bude vaše odpověď ověřena a obdržíte zpětnou vazbu označující, která písmena jsou správná (zelená), která písmena jsou správná, ale na špatné pozici (žlutá), a která písmena jsou nesprávná (červená). Tato zpětná vazba může být užitečná při určování, která písmena zahrnout nebo vyloučit v dalším odhadu.} \\
\\ 
\texttt{Například zpětná vazba pro "jablko" může být:} \\
\texttt{zpětná vazba k hádání: d<žlutá> o<žlutá> l<zelená> a<žlutá> r<červená> o<červená>} \\
\\ 
\texttt{Vysvětlení by mělo obsahovat podrobnosti o tom, jak je zpětná vazba k hádání použita k vytvoření nového odhadu.} \\
\\ 
\texttt{Nevytvářejte žádný jiný text.} \\
\\ 
\texttt{Začněme vaším prvním odhadem.} \\
            }
        }
    }
     \\ \\

    \theutterance \stepcounter{utterance}  
    & \multicolumn{4}{p{0.6\linewidth}}{
        \cellcolor[rgb]{0.561,0.737,0.561}{
            \makecell[{{p{\linewidth}}}]{
                \texttt{\tiny{[P1$\rangle$GM]}}
                \texttt{vysvětlení: Začínám s běžným českým slovem, které obsahuje širokou škálu často se vyskytujících samohlásek a souhlásek pro efektivní eliminaci možností.} \\
\texttt{hádání: cesta} \\
            }
        }
    }
    & & \\ \\

    \theutterance \stepcounter{utterance}  
    & & & \multicolumn{2}{p{0.3\linewidth}}{
        \cellcolor[rgb]{0.9,0.9,0.9}{
            \makecell[{{p{\linewidth}}}]{
                \texttt{\tiny{[GM$|$GM]}}
                \texttt{attempts = 1} \\
\texttt{target\_word = pocit} \\
\texttt{guess = cesta} \\
\texttt{guess\_feedback = c<yellow> e<red> s<red> t<yellow> a<red>} \\
            }
        }
    }
    & & \\ \\

    \theutterance \stepcounter{utterance}  
    & & & \multicolumn{4}{p{0.6\linewidth}}{
        \cellcolor[rgb]{0.9,0.9,0.9}{
            \makecell[{{p{\linewidth}}}]{
                \texttt{\tiny{[P1$\langle$GM]}}
                \texttt{zpětná vazba k hádání: c<žlutá> e<červená> s<červená> t<žlutá> a<červená>} \\
\\ 
\texttt{Odpovězte pouze v tomto formátu.} \\
\texttt{vysvětlení: podrobnosti} \\
\texttt{hádání: slovo} \\
            }
        }
    }
     \\ \\

    \theutterance \stepcounter{utterance}  
    & \multicolumn{4}{p{0.6\linewidth}}{
        \cellcolor[rgb]{0.561,0.737,0.561}{
            \makecell[{{p{\linewidth}}}]{
                \texttt{\tiny{[P1$\rangle$GM]}}
                \texttt{vysvětlení: Písmeno 'c' a 't' jsou v hadnání, ale nejsou na svých pozicích. Písmena 'e', 's' a 'a' jsou vyloučena. Zkusím nové slovo s písmeny 'c' a 't' na jiných pozicích a přidám další časté samohlásky jako 'o' nebo 'i'.} \\
\texttt{hádání: tonic} \\
            }
        }
    }
    & & \\ \\

    \theutterance \stepcounter{utterance}  
    & & & \multicolumn{2}{p{0.3\linewidth}}{
        \cellcolor[rgb]{0.9,0.9,0.9}{
            \makecell[{{p{\linewidth}}}]{
                \texttt{\tiny{[GM$|$GM]}}
                \texttt{Error: The guessed word is not a valid word for this game.} \\
            }
        }
    }
    & & \\ \\

    \theutterance \stepcounter{utterance}  
    & & & \multicolumn{4}{p{0.6\linewidth}}{
        \cellcolor[rgb]{0.9,0.9,0.9}{
            \makecell[{{p{\linewidth}}}]{
                \texttt{\tiny{[P1$\langle$GM]}}
                \texttt{Vaše hádané slovo není pro tuto hru platným slovem. Zkuste to prosím znovu.} \\
\\ 
\texttt{Odpovězte pouze v tomto formátu.} \\
\texttt{vysvětlení: podrobnosti} \\
\texttt{hádání: slovo} \\
            }
        }
    }
     \\ \\

    \theutterance \stepcounter{utterance}  
    & \multicolumn{4}{p{0.6\linewidth}}{
        \cellcolor[rgb]{0.561,0.737,0.561}{
            \makecell[{{p{\linewidth}}}]{
                \texttt{\tiny{[P1$\rangle$GM]}}
                \texttt{vysvětlení: Předchozí slovo nebyl platný český výraz. Musím se vrátit k písmenům 'c' a 't' (která jsou v hadnání, ale na špatné pozici) a vyhnout se písmenům 'e', 's' a 'a'. Zkusím slovo, které využívá písmeno 'o' a 'i' a umístí 'c' a 't' na jiné pozice.} \\
\texttt{hádání: ticho} \\
            }
        }
    }
    & & \\ \\

    \theutterance \stepcounter{utterance}  
    & & & \multicolumn{2}{p{0.3\linewidth}}{
        \cellcolor[rgb]{0.9,0.9,0.9}{
            \makecell[{{p{\linewidth}}}]{
                \texttt{\tiny{[GM$|$GM]}}
                \texttt{attempts = 2} \\
\texttt{target\_word = pocit} \\
\texttt{guess = ticho} \\
\texttt{guess\_feedback = t<yellow> i<yellow> c<green> h<red> o<yellow>} \\
            }
        }
    }
    & & \\ \\

    \theutterance \stepcounter{utterance}  
    & & & \multicolumn{4}{p{0.6\linewidth}}{
        \cellcolor[rgb]{0.9,0.9,0.9}{
            \makecell[{{p{\linewidth}}}]{
                \texttt{\tiny{[P1$\langle$GM]}}
                \texttt{zpětná vazba k hádání: t<žlutá> i<žlutá> c<zelená> h<červená> o<žlutá>} \\
\\ 
\texttt{Odpovězte pouze v tomto formátu.} \\
\texttt{vysvětlení: podrobnosti} \\
\texttt{hádání: slovo} \\
            }
        }
    }
     \\ \\

    \theutterance \stepcounter{utterance}  
    & \multicolumn{4}{p{0.6\linewidth}}{
        \cellcolor[rgb]{0.561,0.737,0.561}{
            \makecell[{{p{\linewidth}}}]{
                \texttt{\tiny{[P1$\rangle$GM]}}
                \texttt{vysvětlení: Písmeno 'c' je na správné pozici (3. místo). Písmena 't', 'i' a 'o' jsou v hadnání, ale na špatných pozicích. Písmeno 'h' je vyloučeno. Musím tedy sestavit slovo, kde bude 'c' na třetím místě, 't', 'i' a 'o' budou na jiných pozicích a vyhnout se 'h', 'e', 's', 'a'. Zkusím slovo, kde 't' bude na začátku, 'i' na konci a 'o' na druhé pozici.} \\
\texttt{hádání: topic} \\
            }
        }
    }
    & & \\ \\

    \theutterance \stepcounter{utterance}  
    & & & \multicolumn{2}{p{0.3\linewidth}}{
        \cellcolor[rgb]{0.9,0.9,0.9}{
            \makecell[{{p{\linewidth}}}]{
                \texttt{\tiny{[GM$|$GM]}}
                \texttt{Error: The guessed word is not a valid word for this game.} \\
            }
        }
    }
    & & \\ \\

    \theutterance \stepcounter{utterance}  
    & & & \multicolumn{4}{p{0.6\linewidth}}{
        \cellcolor[rgb]{0.9,0.9,0.9}{
            \makecell[{{p{\linewidth}}}]{
                \texttt{\tiny{[P1$\langle$GM]}}
                \texttt{Vaše hádané slovo není pro tuto hru platným slovem. Zkuste to prosím znovu.} \\
\\ 
\texttt{Odpovězte pouze v tomto formátu.} \\
\texttt{vysvětlení: podrobnosti} \\
\texttt{hádání: slovo} \\
            }
        }
    }
     \\ \\

    \theutterance \stepcounter{utterance}  
    & \multicolumn{4}{p{0.6\linewidth}}{
        \cellcolor[rgb]{0.561,0.737,0.561}{
            \makecell[{{p{\linewidth}}}]{
                \texttt{\tiny{[P1$\rangle$GM]}}
                \texttt{vysvětlení: Musím se vyhnout cizím slovům a držet se platných českých slov. Z dostupných informací vím: 'c' je na 3. pozici, 't', 'i' a 'o' jsou v hadnání, ale na jiných pozicích. 'h', 'e', 's', 'a' jsou vyloučena. Zkusím slovo, kde 'o' bude na začátku, 't' na konci a 'i' na 2. pozici, s 'c' na 3. pozici.} \\
\texttt{hádání: oticí} \\
            }
        }
    }
    & & \\ \\

    \theutterance \stepcounter{utterance}  
    & & & \multicolumn{2}{p{0.3\linewidth}}{
        \cellcolor[rgb]{0.9,0.9,0.9}{
            \makecell[{{p{\linewidth}}}]{
                \texttt{\tiny{[GM$|$GM]}}
                \texttt{Error: The guessed word is not a valid word for this game.} \\
            }
        }
    }
    & & \\ \\

    \theutterance \stepcounter{utterance}  
    & & & \multicolumn{4}{p{0.6\linewidth}}{
        \cellcolor[rgb]{0.9,0.9,0.9}{
            \makecell[{{p{\linewidth}}}]{
                \texttt{\tiny{[P1$\langle$GM]}}
                \texttt{Vaše hádané slovo není pro tuto hru platným slovem. Zkuste to prosím znovu.} \\
\\ 
\texttt{Odpovězte pouze v tomto formátu.} \\
\texttt{vysvětlení: podrobnosti} \\
\texttt{hádání: slovo} \\
            }
        }
    }
     \\ \\

    \theutterance \stepcounter{utterance}  
    & \multicolumn{4}{p{0.6\linewidth}}{
        \cellcolor[rgb]{0.561,0.737,0.561}{
            \makecell[{{p{\linewidth}}}]{
                \texttt{\tiny{[P1$\rangle$GM]}}
                \texttt{vysvětlení: Musím najít platné české slovo. Známe: 'c' je na 3. pozici, 't', 'i', 'o' jsou v hadnání, ale na jiných pozicích. Vyloučena jsou: 'e', 's', 'a', 'h'. Zkusím slovo, kde 't' je na 1. pozici, 'o' je na 4. pozici a 'i' je na 5. pozici, s 'c' na 3. pozici.} \\
\texttt{hádání: tocí} \\
            }
        }
    }
    & & \\ \\

    \theutterance \stepcounter{utterance}  
    & & & \multicolumn{2}{p{0.3\linewidth}}{
        \cellcolor[rgb]{0.9,0.9,0.9}{
            \makecell[{{p{\linewidth}}}]{
                \texttt{\tiny{[GM$|$GM]}}
                \texttt{Error: The length of the guessed word is not 5.} \\
            }
        }
    }
    & & \\ \\

    \theutterance \stepcounter{utterance}  
    & & & \multicolumn{2}{p{0.3\linewidth}}{
        \cellcolor[rgb]{0.9,0.9,0.9}{
            \makecell[{{p{\linewidth}}}]{
                \texttt{\tiny{[GM$|$GM]}}
                \texttt{game\_result = ABORT} \\
            }
        }
    }
    & & \\ \\

\end{supertabular}
}

\end{document}
